
% ════════════════════════════════════════════
% 引言（01.引言.tex）
% ════════════════════════════════════════════
\section{引言}
\subsection{问题背景}

随着以光伏发电为代表的分布式能源大规模接入配电网，其出力受天气影响的间歇性与随机性给电网的功率平衡带来了严峻挑战。微网作为聚合分布式电源、储能设备与本地负荷的小型发用电系统\cite{lasseter2002}，能够通过本地消纳光伏电量、平抑负荷波动，显著减少对公共电网的依赖与冲击。在非孤岛运行方式下，微网既可以从外部电网购电，也可以将多余的电量反馈电网或存储于储能设备之中，其能量管理与运行调度已成为研究热点\cite{kanchev2011}。

微网与外网调控的核心矛盾在于：一方面，小区负荷与光伏出力在一天之内、一年之内都随时间剧烈变化，光伏只在白天出力且峰谷差异巨大；另一方面，外网电价存在明显的峰谷价差，波动电价下价差更为显著。若微网完全按负荷实时从外网购电，将承担高昂的峰时购电费用；而储能设备的引入，使微网可以在电价较低或光伏有剩余电量的时段充电，在电价较高或光伏不足的时段放电，从而在满足负荷供电的前提下尽可能降低购电费用。然而，储能设备的容量有限、充放电功率有限且存在效率损耗，如何精准制定购电计划与储能充放电计划，使微网在满足约束的同时实现购电费用最小化，是一个亟待解决的优化问题。因此，建立一套科学的微网购电与储能调控模型具有重要的工程意义。

\subsection{问题重述}

\textbf{问题1：}在每日电价与负载相同、光伏预测已知的条件下，建立储能 0:00 与 24:00 储电量相等的微网日计划购电策略模型，并给出指定时段的购电量与储能充放电量。

\textbf{问题2：}在每日电价相同、负载与光伏随时间变化、供电不足按 5 倍电价紧急购电的条件下，建立每天 0:00 制定当天计划购电策略的模型，并给出全年结果。

\textbf{问题3：}在 0:00、6:00、12:00、18:00 可获得未来 24 小时光伏预报的条件下，建立含计划购电、调整购电、违约惩罚与紧急购电的微网购电策略模型，并分析引入其他预报时刻调整的必要性。

\textbf{问题4：}在外网电价实时波动的条件下，建立微网当天购电策略模型，并重新计算问题二与问题三的购电策略。
